\clearpage
\VerdierPhysicalPageAnchor{0189}%
\begingroup\small\setlength{\parskip}{0pt}\setlength{\emergencystretch}{1.5em}
\setlength{\abovedisplayskip}{4pt plus 2pt minus 2pt}%
\setlength{\belowdisplayskip}{4pt plus 2pt minus 2pt}%

\VerdierSemanticAnchor{ch3d-2-1-4-proof-cont}
ここで $v$ はホモトピー同値である。

実際、射 $m$ の写像錐を構成する：

\VerdierSemanticAnchor{ch3d-2-1-4-cone-m}
\[
X\xrightarrow{m}Y\xrightarrow{p(m)}c(m)\longrightarrow X[1],
\]
次に射 $p(m)$ の写像錐を構成する：

\VerdierSemanticAnchor{ch3d-2-1-4-cone-pm}
\[
Y\xrightarrow{p(m)}c(m)\longrightarrow c(p(m))
\xrightarrow{q}Y[1].
\]
射 $c(p(m))\xrightarrow{q}Y[1]$ はエピ射である（第 I 章，3.2.2）。完全三角の
性質（性質 TRII）から、ホモトピー同値

\VerdierSemanticAnchor{ch3d-2-1-4-homotopy-equivalence}
\[
X\xrightarrow{v}c(p(m))[-1]
\]
で、三角

\VerdierSemanticAnchor{ch3d-2-1-4-triangle}
\begin{center}
\begin{tikzcd}[column sep=3.5em,row sep=3.2em]
X \arrow[rr,"v"] \arrow[dr,"m"']
  && c(p(m))[-1] \arrow[dl,"{q[-1]}"] \\
& Y &
\end{tikzcd}
\end{center}
がホモトピーを除いて可換となるものが存在する。

したがって射 $m$ は $\operatorname{comp}(\mathcal A)$ のエピ射であると
仮定してよい。次の複体

\VerdierSemanticAnchor{ch3d-2-1-4-kernel}
\[
N=\cdots\longrightarrow N^i\xrightarrow{d_N^i}N^{i+1}
\xrightarrow{d_N^{i+1}}N^{i+2}\longrightarrow\cdots
\]
を $m$ の核とする。次のようにおく：

\VerdierSemanticAnchor{ch3d-2-1-4-truncation}
\[
N_{-n}=\cdots\longrightarrow N^{n-3}\longrightarrow N^{n-2}
\longrightarrow\operatorname{Im}(d_N^{n-2})
\longrightarrow0\longrightarrow0\longrightarrow\cdots.
\]
$N_{-n}$ は $N$ の非輪状部分複体の減少列をなす。そこで
$X_n=X/N_n$ とおけば、列

\VerdierSemanticAnchor{ch3d-2-1-4-quotients}
\[
u_n\colon X_n\longrightarrow X_{n-1},\qquad n\in\mathbb Z
\]
を得る。各 $u_n$ は高さ $-n$ の初等射であり、

\[
\varprojlim_n X_n=X,
\qquad
\varinjlim_n X_n=Y.
\]
\endgroup
